(1) is Lemma \claimref{transport.regularity-localises-zero-divisors-idempotents}. (2): for $p - q \not\equiv 1 \pmod 4$ the algebra is simple (Theorem \claimref{transport.classification-real-clifford-algebras}), so by Theorem \claimref{transport.existence-classification-type-c} its only central idempotents are $0$ and $1$, excluded in (1). (3): if $a = \varepsilon_- a$ (say), then $\varepsilon_+ a = \varepsilon_+ \varepsilon_- a = 0$ with $\varepsilon_+$ central and nontrivial. If neither central projector annihilates $a$, then both components $\varepsilon_\pm a$ are nonzero; apply Lemma \claimref{transport.regularity-localises-zero-divisors-idempotents} componentwise: in each simple ideal $\varepsilon_\pm A \cong \Mat_k(\mathbb{F})$, the component $\varepsilon_\pm a$ is either invertible (within the ideal) or annihilated by a non-central idempotent of that ideal; since $a$ is a zero divisor of $A$, not both components can be invertible in their ideals, so at least one is annihilated, and the annihilating idempotent, viewed in $A$, is non-central (it is neither $0$ nor $\varepsilon_\pm$ nor central in $A$, central idempotents of $A$ being classified in Theorem \claimref{transport.existence-classification-type-c}: any idempotent of $\varepsilon_\pm A$ central in $A$ would be a central idempotent of $A$ inside one ideal, hence $0$ or $\varepsilon_\pm$; the annihilating idempotent is nontrivial in its ideal, and if it equalled $\varepsilon_\pm$ then $\varepsilon_\pm a$ would be $0$, excluded). The classification of origins (isotropic iff mixed) is Proposition \claimref{transport.existence-locus-type-n} with Example \ref{ex:hyperbolic-idem}.
